% ID: SEL-D0062-OA-E01
% ESTADO: cotejada. Segunda pasada visual y borrador compilado revisados.
% Fuente: Selectividad/sel_2010_matematicas_aplicadas/mat_apli/mat_apli_junio/mat_apli_junio_ra/EXAMEN-ANDALUCÍA-4.pdf
% SHA-256: 4e10e40d60aebdab7570fa77b9ce953d7bf6c9a4c8869a5d3f9b7ca3c907c154
% Página: 1; opción A; ejercicio original 1.

\textbf{(2.5 puntos)} Un comerciante quiere dar salida a 400 kg de avellanas, 300 kg de nueces y 400 kg de almendras. Para ello hace dos tipos de lotes: los de tipo A contienen 2 kg de avellanas, 2 kg de nueces y 1 kg de almendras; y los de tipo B contienen 3 kg de avellanas, 1 kg de nueces y 4 kg de almendras. El precio de venta de cada lote es de 20 euros para los del tipo A y de 40 euros para los del tipo B. \textquestiondown{}Cuántos lotes de cada tipo debe vender para obtener el máximo ingreso y a cuánto asciende éste?
